\corrigesChapitre{3}{Exponentielle imaginaire et trigonométrie}
\begin{corrige}{1}{Passages de formes}
\begin{questions}
\item $1+\mathrm i\sqrt3$.
\item $2\mathrm e^{-\mathrm i\pi/2}$.
\item $\sqrt2\mathrm e^{-\mathrm i\pi/4}$.
\end{questions}
\end{corrige}
\begin{corrige}{2}{Produits et quotients}
\begin{questions}
\item $6\mathrm e^{\mathrm i\pi/2}=6\mathrm i$.
\item $\tfrac23\mathrm e^{-\mathrm i\pi/6}$.
\item $3/2$.
\end{questions}
\end{corrige}
\begin{corrige}{3}{Puissances}
\begin{questions}
\item $\mathrm e^{2\mathrm i\pi}=1$.
\item $16$.
\item $\sqrt3+\mathrm i=2\mathrm e^{\mathrm i\pi/6}$, donc le cube vaut $8\mathrm e^{\mathrm i\pi/2}=8\mathrm i$.
\end{questions}
\end{corrige}
\begin{corrige}{4}{Duplication}
\begin{questions}
\item $2\cos^2x-1$.
\item $2\sin x\cos x$.
\item L’angle vaut $\pi/2$, donc le cosinus vaut $0$ ; les formules d’addition donnent aussi $\sqrt3/4-\sqrt3/4=0$.
\end{questions}
\end{corrige}
\begin{corrige}{5}{Angle $\pi/12$}
\begin{questions}
\item $\pi/4-\pi/6$.
\item $\cos(\pi/12)=(\sqrt6+\sqrt2)/4$.
\item $\sin(\pi/12)=(\sqrt6-\sqrt2)/4$, par la formule de différence.
\end{questions}
\end{corrige}
\begin{corrige}{6}{Euler au carré}
\begin{questions}
\item Développer le carré de $(\mathrm e^{\mathrm ix}+\mathrm e^{-\mathrm ix})/2$ et regrouper les exponentielles opposées.
\item $x/2+\sin2x/4$.
\item $\pi/2$, car le sinus s’annule aux bornes.
\end{questions}
\end{corrige}
\begin{corrige}{7}{Moivre au cube}
\begin{questions}
\item La partie réelle vaut $\cos^3x-3\cos x\sin^2x$ ; l’imaginaire vaut $3\cos^2x\sin x-\sin^3x$.
\item $\cos3x=\cos^3x-3\cos x(1-\cos^2x)=4\cos^3x-3\cos x$.
\item $\cos^3x=(\cos3x+3\cos x)/4$.
\end{questions}
\end{corrige}
\begin{corrige}{8}{Une équation exponentielle}
\begin{questions}
\item $2x\in2\pi\Z$, donc $x=k\pi$.
\item $2x=\pi+2k\pi$, donc $x=\pi/2+k\pi$.
\item $\pi/2$ et $3\pi/2$.
\end{questions}
\end{corrige}
\begin{corrige}{9}{Linéariser une puissance quatre}
\begin{questions}
\item $\cos^4x=(3+4\cos2x+\cos4x)/8$.
\item $3x/8+\sin2x/4+\sin4x/32$.
\item $3\pi/4$ ; les termes sinus sont nuls aux bornes.
\end{questions}
\end{corrige}
\begin{corrige}{10}{Une suite complexe}
\begin{questions}
\item $|z_n|=2^{n/2}$ ; un argument est $n\pi/4$.
\item Lorsque $n\pi/4\in\pi\Z$, soit $n$ multiple de $4$.
\item Lorsque $n\pi/4=\pi/2$ modulo $\pi$, soit $n\equiv2\pmod4$.
\end{questions}
\end{corrige}
\begin{corrige}{11}{Une preuve par le produit scalaire}
\begin{questions}
\item Prendre $(\cos a,\sin a)$ et $(\cos b,\sin b)$.
\item Le calcul par coordonnées vaut $\cos a\cos b+\sin a\sin b$ ; la formule géométrique donne $\cos(a-b)$.
\item Remplacer $b$ par $-b$ et utiliser les parités de cosinus et sinus : $\cos(a+b)=\cos a\cos b-\sin a\sin b$.
\end{questions}
\end{corrige}
\begin{corrige}{12}{Une somme oscillante}
\begin{questions}
\item $q^6=1$ et $q\ne1$ car son argument n’est pas nul modulo $2\pi$.
\item La somme géométrique vaut $(1-q^6)/(1-q)=0$.
\item En prenant les parties réelle et imaginaire, les deux sommes, pour $k=0,\ldots,5$, valent zéro.
\end{questions}
\end{corrige}
